
\section{Limitations}
\label{ssec:limits}
% jp previously said assign attribute e.g color, *size*, but size is controlled by the bbox
%Our method has several limitations at the certain level which inherits all possible shortages from the SD, and the ZeroScope.
Our method shares and inherits  common failure cases of the underlying diffusion model.
Notably, at the time of writing, models based on CLIP and Stable Diffusion sometimes generate deformed objects and struggle to generate multiple objects and correctly assign attributes (e.g.~color) to objects. We show some failures in \FigRef{fig:limit}. For instance, we requested a red jeep driving on the road but the synthesis shows it sinking into a mud road. The panda example shows the camera moving instead of the panda itself. The red car has implausible deformation, and Darth Vader's light saber turns into a surf board. The length of the resulting video clips is restricted to that produced by the pre-trained model, for instance, the 24 images in the case of ZeroScope. This is not a crucial limitation, as movies are commonly (with some exceptions!) composed of short ``shots'' of several seconds each.  The bbox guides object placement without precisely constraining it. This is an advantage as well, however, since otherwise the user would have to specify the correct x-y aspect ratio for objects, a complicated task for non-artists.


\begin{figure}[tp]
  \centering
  % \includegraphics[width=\linewidth]{assets/limitation.jpg }
  \includegraphics[width=\linewidth]{assets/limitation.pdf}
  \vspace{-0.5em}
  \caption{\textbf{Failure cases.} %Several limitations depending on the context.
  Prompts used in subfigures: ``\textbf{A red jeep} driving on the road'', ``\textbf{A red car} driving on the highway'', ``\textbf{a panda} eating bamboo'', and ``\textbf{Darth Vader} surfing in waves'', where the bold prompt word is the directed subject.}
  \label{fig:limit}
\vspace{-1em}
\end{figure}

